\section*{Your turn}

Check out tag \code{ch15} of the companion repository and run \code{npm ci}
once. Exercises 1 to 7 need the database, the seven services and the router.
Exercises 8 and 9 need node and nothing started at all.

Everything in this chapter is signed with one key, and nothing will accept a
token until that key is in your environment:

\begin{minted}[fontsize=\footnotesize,breaklines]{text}
export MOSAIC_JWT_SECRET=dev-secret-not-for-anything-real-0123456789
docker compose up -d --build mosaic-db mosaic-nats mosaic-catalog \
    mosaic-pricing mosaic-inventory mosaic-accounts mosaic-reviews \
    mosaic-ordering mosaic-nodes mosaic-router
docker compose ps
\end{minted}

If the seven services will not start at all, that variable is missing from
their environment: each of them says so by name and then stops. If they start
and every token you mint is refused with HTTP 401, the variable is set in both
places and the two values differ.

\begin{enumerate}
  \item \textbf{Get all four answers.} Send the catalogue query from
  section~\ref{sec:ch15-four-callers-one-query} four times: with no header,
  with \code{--scope "read:pii"}, with no scope, and with
  \code{--expires-in -60}. Before you look at the bodies, write down which HTTP
  status you expect for each. Two of the four surprise most people.

  \item \textbf{Find your own blast radius.} The anonymous run loses
  \code{author}, not just \code{email}, because the null had nowhere to stop
  until it reached a nullable field. Make \code{Review.author} non-null in
  \code{schema/reviews.graphql}, recompose, restart the router, and run the
  anonymous query again. Then put it back. What you have just done is decision
  70's argument in reverse, and the result is why that decision matters more
  than it looked at the time.

  \item \textbf{Break the signing key on purpose.} In
  \code{router/config.yaml}, change the jwks entry's \code{secret} from the
  shell form to the Go template form that
  section~\ref{sec:ch15-a-token-the-router-accepts} shows, restart the router,
  and try
  a token that worked a minute ago. Now go looking for the cause using only the
  router's logs at \code{debug}. Give up when you want to. This is the failure
  from section~\ref{sec:ch15-a-token-the-router-accepts} and the point of the
  exercise is how little there is to find.

  \item \textbf{Take the header rule away.} Comment out the \code{headers}
  block, restart the router, and ask for \code{customerById} with a good token.
  Read the error a client would get and decide how long it would take you to
  reach \enquote{the router is not forwarding a header} from those words alone.

  \item \textbf{Move the rule to the other side.} Take \code{[Authorize]} off
  \code{Query.customerById} and put \code{[Authenticated]} from
  \code{HotChocolate.ApolloFederation} on it instead. Export the schema,
  recompose, and diff the composed client schema. The field is now protected by
  the router rather than by Accounts. Ask Accounts directly on 5104 with no
  token and see what you can still read.

  \item \textbf{Turn the graph internal.} Set
  \code{authorization.require\_authentication: true} and restart. Every
  anonymous request now fails, including the storefront query the whole book
  has been running. Decide, in writing, which of Mosaic's fields you would
  actually want behind that switch, and whether the answer is a switch at all.

  \item \textbf{Empty the scopes list.} Run one case and read what it asserts:

  \begin{minted}[fontsize=\footnotesize,breaklines]{text}
node scripts/auth-cases.mjs --print empty-scopes-composes-and-vanishes
  \end{minted}

  Then do it by hand: put \code{scopes: []} on \code{Customer.email},
  compose, and search the output for the directive's name. Nothing in the
  toolchain will tell you the field is public again.

  \item \textbf{Verify it in Postman.} The collection for this chapter is
  \code{postman/mosaic-auth}, six requests against the router. Three of them
  need a token, and the collection cannot mint one, because signing needs the
  key and a collection is not where a key belongs:

  \begin{minted}[fontsize=\footnotesize,breaklines]{text}
npx newman run postman/mosaic-auth.postman_collection.json \
    -e postman/mosaic-auth.local.postman_environment.json \
    --env-var "scopedToken=$(node scripts/mint-token.mjs --scope 'read:pii')" \
    --env-var "plainToken=$(node scripts/mint-token.mjs)" \
    --env-var "expiredToken=$(node scripts/mint-token.mjs --expires-in -60)"
  \end{minted}

  All green at \code{ch15}. Now comment out the \code{headers} block, restart
  the router, and run it again: exactly one request fails, and its assertions
  say more about why than the error a client would see.

  \item \textbf{Guess the merge, then check.} Two subgraphs guard one shared
  field with one scope each. Write down whether a caller should need both or
  either. Then run the same script again and read the two arrays it asserts:

  \begin{minted}[fontsize=\footnotesize,breaklines]{text}
node scripts/auth-cases.mjs --print partial-requires-scopes-merges-to-and
  \end{minted}

  If you guessed the one named \code{...ByOr}, you guessed the field the router
  never reads.
\end{enumerate}

If exercise 1 gives you four identical answers, no token is reaching the
router: the header value has to begin with the word \code{Bearer}, which is
what \code{header\_value\_prefix} expects, and a bare token is discarded
without a word. If exercise 5 composes but Accounts still refuses the field,
\code{[Authorize]} is still on it: the two attributes are independent and
removing one does not disturb the other, which is this chapter's whole argument
arriving as a thing you have to check twice.

The thread this chapter does not close is \code{Query.node}, and the reason is
in section~\ref{sec:ch15-what-the-composer-checks}. The other one is a
subscription: chapter~\ref{ch:real-time-in-a-federated-world} left a connection
that outlives the token that opened it, and nothing here provoked it.
Chapter~\ref{ch:resilience-and-performance} owns what happens when a token
expires with a socket still open.
